\begin{conclusion}[统一伸缩下 Pick--Poisson Gram 谱的严格同伦缩放]\label{con:xi-terminal-pick-poisson-homothetic-spectrum-scaling}
沿用结论 \ref{con:xi-terminal-visibility-complexity-logm-translation} 的统一伸缩 \(\phi_m\) 与 Gram 矩阵 \(\mathsf P(m)\) 的记号，并令 \(\chi_{\mathsf P}(\lambda):=\det(\lambda I_\kappa-\mathsf P)\)。
则特征多项式满足严格协变律
\[
\boxed{\;
\chi_{\mathsf P(m)}(\lambda)=m^{-\kappa}\chi_{\mathsf P}(m\lambda)\; } .
\]
从而按重数计有谱严格缩放
\[
\boxed{\;\mathrm{Spec}(\mathsf P(m))=m^{-1}\mathrm{Spec}(\mathsf P).\;}
\]
\end{conclusion}
\begin{proof}
由命题 \ref{prop:xi-pick-poisson-det-scaling-law}，\(\rho\) 在 \(\phi_m\) 下不变且端点 Poisson 权重满足 \(p_j(m)=m^{-1}p_j\)。
因此对任意指标子集 \(I\subset[\kappa]\)，命题 \ref{prop:xi-pick-poisson-principal-minors-partition} 的闭式给出
\(
\det(\mathsf P_I(m))=m^{-|I|}\det(\mathsf P_I)
\)。
令 \(Z_r\) 与 \(Z_r(m)\) 为命题 \ref{prop:xi-pick-poisson-charpoly-coefficients} 中的 \(r\)-阶主子式配分函数，则
\(
Z_r(m)=m^{-r}Z_r
\)。
代入 \(\chi_{\mathsf P}(\lambda)=\sum_{r=0}^\kappa (-1)^r Z_r \lambda^{\kappa-r}\) 得
\[
\chi_{\mathsf P(m)}(\lambda)
=\sum_{r=0}^\kappa (-1)^r m^{-r}Z_r\,\lambda^{\kappa-r}
=m^{-\kappa}\sum_{r=0}^\kappa(-1)^r Z_r\,(m\lambda)^{\kappa-r}
=m^{-\kappa}\chi_{\mathsf P}(m\lambda).
\]
根随变量按 \(m^{-1}\) 缩放，遂得谱缩放结论。证毕。
\end{proof}

\begin{conclusion}[全部谱函数的重整化恒等式]\label{con:xi-terminal-pick-poisson-spectral-function-renormalization}
在结论 \ref{con:xi-terminal-pick-poisson-homothetic-spectrum-scaling} 的设定下，对任意在 \(\mathrm{Spec}(\mathsf P)\) 的邻域内解析的函数 \(f\)，有
\[
\boxed{\;
\mathrm{tr}\,f(\mathsf P(m))=\mathrm{tr}\,f(m^{-1}\mathsf P)\;}
\]
并且对任意复参数 \(z\) 有
\[
\boxed{\;
\det(I+z\mathsf P(m))=\det\!\bigl(I+(z/m)\mathsf P\bigr)\; } .
\]
\end{conclusion}
\begin{proof}
由谱缩放，存在重排使 \(\lambda_j(\mathsf P(m))=m^{-1}\lambda_j(\mathsf P)\)。
于是 \(\mathrm{tr}\,f(\mathsf P(m))=\sum_j f(m^{-1}\lambda_j(\mathsf P))=\mathrm{tr}\,f(m^{-1}\mathsf P)\)。
行列式恒等式同理由 \(\det(I+z\mathsf P(m))=\prod_j(1+z\lambda_j(\mathsf P(m)))\) 得到。证毕。
\end{proof}

\begin{conclusion}[尺度不变量与平移型势的规范唯一性]\label{con:xi-terminal-pick-poisson-scale-invariant-and-affine-uniqueness}
沿用结论 \ref{con:xi-terminal-visibility-complexity-logm-translation} 中的
\(
S_{\partial}=-\log\det\mathsf P
\)、
\(
V_{\partial}=\log\mathrm{tr}(\mathsf P)
\)
与 \(\phi_m\) 的记号。
定义尺度不变量
\[
\boxed{\;
I_{\partial}:=S_{\partial}+\kappa V_{\partial}
=-\log\!\left(\frac{\det\mathsf P}{(\mathrm{tr}\,\mathsf P)^{\kappa}}\right)\; } .
\]
则对一切 \(m>0\) 有严格不变性 \(I_{\partial}(m)\equiv I_{\partial}(1)\)。
\par
更一般地，若线性组合 \(F:=aS_{\partial}+bV_{\partial}\) 满足对任意配置与任意 \(m>0\) 皆有
\[
F(m)=F(1)+c\log m\qquad(c\ \text{与配置无关}),
\]
则必有 \(c=a\kappa-b\)。
因此：
\begin{enumerate}
  \item \(c=0\) 当且仅当 \((a,b)\) 与 \((1,\kappa)\) 成比例（即 \(F\) 与 \(I_{\partial}\) 成比例）；
  \item \(c=a\) 当且仅当 \((a,b)\) 与 \((1,\kappa-1)\) 成比例（即 \(F\) 与 \(\mathscr C_{\partial}=S_{\partial}+(\kappa-1)V_{\partial}\) 成比例）。
\end{enumerate}
\end{conclusion}
\begin{proof}
由命题 \ref{prop:xi-pick-poisson-det-scaling-law} 得
\(
S_{\partial}(m)=S_{\partial}(1)+\kappa\log m
\)、
\(
V_{\partial}(m)=V_{\partial}(1)-\log m
\)。
故
\(
I_{\partial}(m)=I_{\partial}(1)
\)
且
\(
F(m)=F(1)+(a\kappa-b)\log m
\)，从而 \(c=a\kappa-b\)。
两条“当且仅当”分别对应 \(a\kappa-b=0\) 与 \(a\kappa-b=a\)。证毕。
\end{proof}

\begin{conclusion}[特征行列式的齐次性与一阶重整化方程]\label{con:xi-terminal-pick-poisson-characteristic-homogeneity-pde}
设
\[
\Psi(z,m):=\det(I-z\mathsf P(m)).
\]
则有严格齐次性
\[
\boxed{\;\Psi(z,m)=\Psi(z/m,1)\;}
\]
并因此满足一阶重整化方程
\[
\boxed{\;
\bigl(m\partial_m+z\partial_z\bigr)\Psi(z,m)=0,\qquad
\bigl(m\partial_m+z\partial_z\bigr)\log\Psi(z,m)=0\; } .
\]
\end{conclusion}
\begin{proof}
由结论 \ref{con:xi-terminal-pick-poisson-spectral-function-renormalization} 取 \(f(x)=\log(1-zx)\)（或直接对谱乘积）得
\(
\det(I-z\mathsf P(m))=\det(I-(z/m)\mathsf P)
\)，即 \(\Psi(z,m)=\Psi(z/m,1)\)。
对该恒等式对 \(m,z\) 求导并消去 \(\partial_{(\cdot)}\Psi(z/m,1)\) 即得所示偏微分方程。证毕。
\end{proof}

\begin{corollary}[零点集的严格缩放]\label{cor:xi-pick-poisson-characteristic-zero-set-scaling}
在结论 \ref{con:xi-terminal-pick-poisson-characteristic-homogeneity-pde} 的设定下，定义零点集
\[
Z(m):=\{z\in\CC:\ \Psi(z,m)=0\}.
\]
则对一切 \(m>0\) 有严格比例律
\[
Z(m)=m\cdot Z(1).
\]
\end{corollary}
\begin{proof}
由齐次性 \(\Psi(z,m)=\Psi(z/m,1)\)，有
\[
z\in Z(m)\ \Longleftrightarrow\ \Psi(z,m)=0
\Longleftrightarrow\ \Psi(z/m,1)=0
\Longleftrightarrow\ z/m\in Z(1),
\]
故 \(Z(m)=mZ(1)\)。证毕。
\end{proof}

\begin{corollary}[谱倒数的对数平移不变量]\label{cor:xi-pick-poisson-spectrum-log-translation-invariant}
在上述设定下，令 \(\{\lambda_j(m)\}_{j=1}^{\kappa}=\mathrm{Spec}(\mathsf P(m))\)（计重数）。
则按重数计有
\[
\Bigl\{\frac{1}{\lambda_j(m)}\Bigr\}_{j=1}^{\kappa}=Z(m)=mZ(1)
=\Bigl\{\frac{m}{\lambda_j(1)}\Bigr\}_{j=1}^{\kappa},
\]
从而
\[
\lambda_j(m)=\frac{1}{m}\lambda_j(1)\qquad(1\le j\le\kappa),
\]
并且谱的对数计数测度
\[
\mu_m:=\frac{1}{\kappa}\sum_{j=1}^{\kappa}\delta_{\log \lambda_j(m)}
\]
满足推送恒等式
\[
\mu_m=\Bigl(x\mapsto x-\log m\Bigr)_{\#}\mu_1.
\]
\end{corollary}
\begin{proof}
由 \(\Psi(z,m)=\det(I-z\mathsf P(m))=\prod_{j=1}^{\kappa}(1-z\lambda_j(m))\)，其零点（计重数）恰为 \(\{1/\lambda_j(m)\}_{j=1}^{\kappa}\)。
结合推论 \ref{cor:xi-pick-poisson-characteristic-zero-set-scaling} 得 \(\{1/\lambda_j(m)\}=m\{1/\lambda_j(1)\}\)，逐项取倒数即得 \(\lambda_j(m)=m^{-1}\lambda_j(1)\)；
对数推送公式由 \(\log\lambda_j(m)=\log\lambda_j(1)-\log m\) 直接推出。证毕。
\end{proof}

\begin{conclusion}[Schur 补账本的逐项同伦缩放]\label{con:xi-terminal-pick-poisson-schur-ledger-homothetic-scaling}
沿用结论 \ref{con:xi-terminal-pick-poisson-schur-ledger} 中的 Schur 补通量 \(\pi_m\)。
在统一伸缩 \(\phi_m\) 下，用同一顺序定义 \(\pi_r(m)\)。
则对每个 \(r=1,\dots,\kappa\) 有严格缩放律
\[
\boxed{\;\pi_r(m)=m^{-1}\pi_r(1),\qquad -\log\pi_r(m)=-\log\pi_r(1)+\log m.\;}
\]
\end{conclusion}
\begin{proof}
由 Schur 补链式分解，
\(
\pi_r=\det\mathsf P^{(r)}/\det\mathsf P^{(r-1)}
\)
（其中 \(\mathsf P^{(r)}\) 为前 \(r\) 个点的主子矩阵）。
而命题 \ref{prop:xi-pick-poisson-det-scaling-law} 的协变律同样适用于任意主子矩阵，给出
\(
\det\mathsf P^{(r)}(m)=m^{-r}\det\mathsf P^{(r)}
\)。
代入比值即得 \(\pi_r(m)=m^{-1}\pi_r\)，取对数得到第二式。证毕。
\end{proof}

\begin{conclusion}[尺度中心化账本的严格不变量分解]\label{con:xi-terminal-pick-poisson-centered-ledger-invariants}
定义尺度中心化通量增量
\[
\eta_r:=-\log\pi_r+V_{\partial}\qquad(1\le r\le\kappa).
\]
则对一切 \(m>0\) 与一切 \(r\) 有严格尺度不变性 \(\eta_r(m)=\eta_r(1)\)。
并且满足严格求和分解
\[
\boxed{\;
\sum_{r=1}^{\kappa}\eta_r
=\sum_{r=1}^{\kappa}(-\log\pi_r)+\kappa V_{\partial}
=S_{\partial}+\kappa V_{\partial}
=I_{\partial}\; } .
\]
\end{conclusion}
\begin{proof}
由结论 \ref{con:xi-terminal-pick-poisson-schur-ledger-homothetic-scaling}，
\(
-\log\pi_r(m)=-\log\pi_r(1)+\log m
\)；
由命题 \ref{prop:xi-pick-poisson-det-scaling-law}，
\(
V_{\partial}(m)=V_{\partial}(1)-\log m
\)。
两式相加即得 \(\eta_r(m)=\eta_r(1)\)。
求和恒等式由 \(\det\mathsf P=\prod_r\pi_r\) 与 \(I_{\partial}=S_{\partial}+\kappa V_{\partial}\) 直接推出。证毕。
\end{proof}

\begin{conclusion}[尺度不变量的谱型 KL 表示与等号刻画]\label{con:xi-terminal-pick-poisson-scale-invariant-spectral-kl}
设 \(\mathsf P\succ 0\) 的特征值为 \(\lambda_1,\dots,\lambda_\kappa>0\)，并令
\[
p_\Sigma:=\mathrm{tr}(\mathsf P)=\sum_{j=1}^{\kappa}\lambda_j,
\qquad
\mu_j:=\frac{\lambda_j}{p_\Sigma}\ \ (1\le j\le\kappa).
\]
记 \(U_\kappa\) 为 \(\{1,\dots,\kappa\}\) 上的均匀分布，则尺度不变量满足严格恒等式
\[
\boxed{\;
I_{\partial}
\;=\;
\kappa\log\kappa+\kappa\,D_{\mathrm{KL}}(U_\kappa\Vert \mu)\; } ,
\]
其中
\[
D_{\mathrm{KL}}(U_\kappa\Vert \mu)
:=\sum_{j=1}^{\kappa}\frac1\kappa\log\frac{1/\kappa}{\mu_j}\ \ge\ 0.
\]
从而得到谱不等式
\[
\boxed{\;
\det\mathsf P\ \le\ \left(\frac{\mathrm{tr}(\mathsf P)}{\kappa}\right)^{\kappa}\; } ,
\]
且等号当且仅当 \(\mu_1=\cdots=\mu_\kappa=\frac1\kappa\)，亦即
\[
\mathsf P=\frac{\mathrm{tr}(\mathsf P)}{\kappa}\,I_\kappa .
\]
\end{conclusion}
\begin{proof}
由 \(\det\mathsf P=\prod_{j=1}^{\kappa}\lambda_j=(p_\Sigma)^\kappa\prod_{j=1}^{\kappa}\mu_j\) 得
\[
I_{\partial}
:=-\log\!\left(\frac{\det\mathsf P}{(\mathrm{tr}\,\mathsf P)^{\kappa}}\right)
=-\sum_{j=1}^{\kappa}\log\mu_j.
\]
另一方面，
\[
D_{\mathrm{KL}}(U_\kappa\Vert \mu)
=\sum_{j=1}^{\kappa}\frac1\kappa\log\frac{1/\kappa}{\mu_j}
=-\log\kappa-\frac1\kappa\sum_{j=1}^{\kappa}\log\mu_j,
\]
故 \(-\sum_j\log\mu_j=\kappa\log\kappa+\kappa D_{\mathrm{KL}}(U_\kappa\Vert \mu)\)，代回即得所示恒等式。
由 \(D_{\mathrm{KL}}\ge 0\) 得 \(-\sum_j\log\mu_j\ge \kappa\log\kappa\)，等价于 \(\prod_j\mu_j\le \kappa^{-\kappa}\)，从而得到 \(\det\mathsf P\le (\mathrm{tr}(\mathsf P)/\kappa)^\kappa\)。
等号条件 \(D_{\mathrm{KL}}=0\) 等价于 \(\mu\) 为均匀分布，亦即 \(\mathsf P\) 的谱完全平坦，从而 \(\mathsf P=(\mathrm{tr}(\mathsf P)/\kappa)I_\kappa\)。证毕。
\end{proof}

\begin{conclusion}[显影复杂度的 KL--互能--标度三项分解]\label{con:xi-terminal-pick-poisson-visibility-complexity-kl-energy-split}
记端点 Poisson 权重
\(
p_j:=\mathsf P_{jj}=P_{a_j}(-1)
\)
并令 \(p_\Sigma=\sum_{j=1}^{\kappa}p_j=\mathrm{tr}(\mathsf P)\)，再定义归一化权重分布
\[
q_j:=\frac{p_j}{p_\Sigma}\qquad(1\le j\le\kappa).
\]
则显影复杂度
\(
\mathscr{C}_{\partial}=S_{\partial}+(\kappa-1)V_{\partial}
\)
满足严格恒等式
\[
\boxed{\;
\mathscr{C}_{\partial}
=\kappa\log\kappa-\log p_\Sigma
\;+\;\kappa\,D_{\mathrm{KL}}(U_\kappa\Vert q)
\;+\;2\sum_{1\le j<k\le\kappa}\bigl(-\log\rho(a_j,a_k)\bigr)\; } .
\]
其中第一项为纯标度项，第二项为端点权重的均匀性偏差，第三项为伪双曲互能。
\end{conclusion}
\begin{proof}
由结论 \ref{con:xi-terminal-negative-spectrum-det-invariant} 的行列式分解
\[
\det\mathsf P=\Bigl(\prod_{j=1}^{\kappa}p_j\Bigr)\Bigl(\prod_{j<k}\rho(a_j,a_k)^2\Bigr)
\]
得
\[
S_{\partial}=-\log\det\mathsf P
=-\sum_{j=1}^{\kappa}\log p_j
2\sum_{j<k}\bigl(-\log\rho(a_j,a_k)\bigr).
\]
于是
\[
\mathscr{C}_{\partial}
=S_{\partial}+(\kappa-1)\log p_\Sigma
=-\sum_{j=1}^{\kappa}\log p_j+(\kappa-1)\log p_\Sigma
2\sum_{j<k}\bigl(-\log\rho(a_j,a_k)\bigr).
\]
又 \(-\sum_j\log p_j=-\kappa\log p_\Sigma-\sum_j\log q_j\)，故
\[
\mathscr{C}_{\partial}
=-\log p_\Sigma-\sum_{j=1}^{\kappa}\log q_j
2\sum_{j<k}\bigl(-\log\rho(a_j,a_k)\bigr).
\]
最后
\[
D_{\mathrm{KL}}(U_\kappa\Vert q)
=-\log\kappa-\frac1\kappa\sum_{j=1}^{\kappa}\log q_j
\]
给出 \(-\sum_j\log q_j=\kappa\log\kappa+\kappa D_{\mathrm{KL}}(U_\kappa\Vert q)\)，代回即得结论。证毕。
\end{proof}

\begin{conclusion}[Schur 外部化的归一化乘积恒等式与有限步塌缩见证]\label{con:xi-terminal-pick-poisson-schur-vandermonde-externalization}
对任意排列 \(\sigma\in\mathfrak S_\kappa\)，令 \(\mathsf P^{(m)}_\sigma\) 为按顺序 \(\sigma(1),\dots,\sigma(m)\) 取出的主子矩阵，并定义相应 Schur 补通量
\[
\pi_m(\sigma)
:=\det(\mathsf P^{(m)}_\sigma)\Big/\det(\mathsf P^{(m-1)}_\sigma)
\qquad(1\le m\le\kappa),
\]
其中约定 \(\det(\mathsf P^{(0)}_\sigma)=1\)。
则满足严格归一化恒等式
\[
\boxed{\;
\prod_{m=1}^{\kappa}\frac{\pi_m(\sigma)}{p_{\sigma(m)}}
=\prod_{1\le j<k\le\kappa}\rho(a_j,a_k)^2\; } .
\]
\end{conclusion}
\begin{proof}
由 Schur 补链式分解可得 \(\det\mathsf P=\prod_{m=1}^{\kappa}\pi_m(\sigma)\)（对置换后的矩阵施行同一分解）。
另一方面，结论 \ref{con:xi-terminal-negative-spectrum-det-invariant} 给出
\(
\det\mathsf P=(\prod_j p_j)(\prod_{j<k}\rho(a_j,a_k)^2)
\)。
将两式相除并注意 \(\prod_m p_{\sigma(m)}=\prod_j p_j\)，得到归一化乘积恒等式。
\par
证毕。
\end{proof}

\begin{proposition}[Schur 消元熵的排列不变性与瓶颈步界]\label{prop:xi-schur-elimination-entropy-permutation-invariant-bottleneck}
\leanverified{paper\_xi\_schur\_elimination\_entropy\_permutation\_invariant\_bottleneck}
沿用结论 \ref{con:xi-terminal-pick-poisson-schur-vandermonde-externalization} 的设定与记号。
对任意排列 \(\sigma\in\mathfrak S_\kappa\)，定义 \textbf{Schur 消元熵}
\[
\mathcal E_{\mathrm{elim}}
:=\frac1\kappa\sum_{m=1}^{\kappa}\log\frac{p_{\sigma(m)}}{\pi_m(\sigma)}.
\]
则：
\begin{enumerate}
  \item \(\mathcal E_{\mathrm{elim}}\) 与排列 \(\sigma\) 无关，并满足闭式
  \[
  \mathcal E_{\mathrm{elim}}
  =\frac{-1}{\kappa}\log\!\Bigl(\prod_{1\le j<k\le\kappa}\rho(a_j,a_k)^2\Bigr)
  =\frac{2}{\kappa}\sum_{1\le j<k\le\kappa}\log\frac{1}{\rho(a_j,a_k)}.
  \]
  \item 对任意排列 \(\sigma\)，必存在某一步 \(m_\star\in\{1,\dots,\kappa\}\) 使
  \[
  \frac{\pi_{m_\star}(\sigma)}{p_{\sigma(m_\star)}}
  \ \le\
  \exp\!\bigl(-\mathcal E_{\mathrm{elim}}\bigr)
  =\Bigl(\prod_{1\le j<k\le\kappa}\rho(a_j,a_k)^2\Bigr)^{1/\kappa}.
  \]
\end{enumerate}
\end{proposition}
\begin{proof}
由结论 \ref{con:xi-terminal-pick-poisson-schur-vandermonde-externalization} 的乘积恒等式，
\(
\prod_{m=1}^{\kappa}\frac{\pi_m(\sigma)}{p_{\sigma(m)}}=\prod_{j<k}\rho(a_j,a_k)^2
\)；
取对数并除以 \(\kappa\) 得
\[
\frac1\kappa\sum_{m=1}^{\kappa}\log\frac{p_{\sigma(m)}}{\pi_m(\sigma)}
=\frac{-1}{\kappa}\log\!\Bigl(\prod_{j<k}\rho(a_j,a_k)^2\Bigr)
=\frac{2}{\kappa}\sum_{j<k}\log\frac1{\rho(a_j,a_k)},
\]
从而 \(\mathcal E_{\mathrm{elim}}\) 与 \(\sigma\) 无关并满足闭式。
\par
对数列 \(\{\log(p_{\sigma(m)}/\pi_m(\sigma))\}_{m=1}^{\kappa}\) 取平均，存在某个 \(m_\star\) 使该项不小于平均值 \(\mathcal E_{\mathrm{elim}}\)；
指数化即得所示瓶颈步界。证毕。
\end{proof}

\begin{conclusion}[Toeplitz 负谱几何均值的严格同伦缩放]\label{con:xi-terminal-toeplitz-negative-geometric-mean-homothety}
设 Toeplitz 主子矩阵 \(\mathbf{T}_N\) 进入稳定显影区间且其负特征值（按重数）为 \(\{\lambda_j^{-}(\mathbf{T}_N)\}_{j=1}^{\kappa}\)，则
\[
\lim_{N\to\infty}\ \prod_{j=1}^{\kappa}\bigl(-\lambda_j^{-}(\mathbf{T}_N)\bigr)=\det\mathsf P.
\]
对统一伸缩后的缺陷数据，记相应 Toeplitz 系列为 \(\mathbf{T}_N(m)\)，其负特征值为 \(\{\lambda_j^{-}(\mathbf{T}_N(m))\}_{j=1}^{\kappa}\)。
则负谱几何均值满足严格缩放律
\[
\boxed{\;
\lim_{N\to\infty}\left(\prod_{j=1}^{\kappa}\bigl(-\lambda_j^{-}(\mathbf{T}_N(m))\bigr)\right)^{1/\kappa}
=m^{-1}\,
\lim_{N\to\infty}\left(\prod_{j=1}^{\kappa}\bigl(-\lambda_j^{-}(\mathbf{T}_N)\bigr)\right)^{1/\kappa}\; } .
\]
\end{conclusion}
\begin{proof}
对伸缩后的系统同样有乘积极限，其极限为 \(\det\mathsf P(m)\)（命题 \ref{prop:xi-pick-poisson-det-factorization} 的逐配置应用）。
而命题 \ref{prop:xi-pick-poisson-det-scaling-law} 给出 \(\det\mathsf P(m)=m^{-\kappa}\det\mathsf P\)。
故乘积极限相差因子 \(m^{-\kappa}\)，两边取 \(1/\kappa\) 次方即得结论。证毕。
\end{proof}

\begin{proposition}[簇分离下交叉互能的显式指数上界]\label{prop:xi-pick-poisson-cross-energy-exp-bound}
设缺陷多重集分解为不交并 \(A\sqcup B\)。
令双曲距离
\[
d_{\mathbb{D}}(a,b):=2\,\operatorname{artanh}\rho(a,b)\qquad(a,b\in\mathbb D),
\]
并设簇间最小分离度
\[
L:=\inf\{d_{\mathbb D}(a,b):a\in A,\ b\in B\}.
\]
则当 \(L\ge \log 2\) 时有显式上界
\[
\boxed{\;
0\le -\log\det\mathsf P(A\cup B)+\log\det\mathsf P(A)+\log\det\mathsf P(B)
\le 8|A||B|\,e^{-L}\; } .
\]
\end{proposition}
\begin{proof}
由命题 \ref{prop:xi-pick-poisson-union-submultiplicativity-cross-energy}，
交叉缺口等于
\(
2\sum_{a\in A}\sum_{b\in B}(-\log\rho(a,b))
\)。
由 \(\rho=\tanh(d_{\mathbb D}/2)\)，令 \(q:=e^{-d_{\mathbb D}}\in(0,1)\)，则
\[
-\log\rho=-\log\tanh\!\Bigl(\frac{d_{\mathbb D}}{2}\Bigr)=\log\frac{1+q}{1-q}.
\]
若 \(d_{\mathbb D}\ge L\ge\log 2\)，则 \(q\le e^{-\log2}=1/2\)，从而
\[
\log\frac{1+q}{1-q}\le q+\frac{q}{1-q}\le 4q\le 4e^{-L}.
\]
代回求和得到
\(
2\sum_{a,b}(-\log\rho(a,b))\le 2|A||B|\cdot 4e^{-L}
\)，即得结论。证毕。
\end{proof}

\begin{corollary}[簇分离下尺度不变量的近可加性]\label{cor:xi-pick-poisson-scale-invariant-near-additivity}
在命题 \ref{prop:xi-pick-poisson-cross-energy-exp-bound} 的条件下，尺度不变量
\(
I_{\partial}(X)=S_{\partial}(X)+|X|\,V_{\partial}(X)
\)
满足估计
\[
\boxed{\;
\bigl|I_{\partial}(A\cup B)-I_{\partial}(A)-I_{\partial}(B)\bigr|
\le 8|A||B|\,e^{-L}
\;+\;(|A|+|B|)\left|\log\frac{\mathrm{tr}\,\mathsf P(A)\,\mathrm{tr}\,\mathsf P(B)}{\mathrm{tr}\,\mathsf P(A\cup B)}\right|\; } .
\]
特别地，若进一步满足 \(\mathrm{tr}\,\mathsf P(A\cup B)=\mathrm{tr}\,\mathsf P(A)+\mathrm{tr}\,\mathsf P(B)\)，则右端第二项为零，\(I_{\partial}\) 的偏离仅由 \(e^{-L}\) 控制。
\end{corollary}
\begin{proof}
由定义
\(
I_{\partial}=S_{\partial}+\kappa V_{\partial}
\)
且 \(|A\cup B|=|A|+|B|\)，有恒等分解
\[
I_{\partial}(A\cup B)-I_{\partial}(A)-I_{\partial}(B)
=\bigl[S_{\partial}(A\cup B)-S_{\partial}(A)-S_{\partial}(B)\bigr]
\]
\[
\qquad
+(|A|+|B|)\log\mathrm{tr}\,\mathsf P(A\cup B)-|A|\log\mathrm{tr}\,\mathsf P(A)-|B|\log\mathrm{tr}\,\mathsf P(B).
\]
第一括号由命题 \ref{prop:xi-pick-poisson-cross-energy-exp-bound} 上界控制；
第二项取绝对值并提取 \((|A|+|B|)\) 即得所示估计。证毕。
\end{proof}

\endinput
